\section{The Shadow Enhanced Greedy Quantum Eigensolver}
\label{sec:main}
\subsection{Procedure}
\label{sec:procedure}
Starting from an empty circuit $C_0=I$, the Shadow Enhanced Greedy Quantum Eigensolver (SEGQE) iteratively constructs an approximation to the ground state of a given $n$-qubit Hamiltonian $H=\sum^r_{i=1}c_iP_i$ by growing a quantum circuit $C$. We denote by $C_k$ the circuit after $k$ iterations. Applied to an initial state $\ket{\psi_0}$, which may encode prior classical information such as a Hartree--Fock state or a matrix product state, the circuit prepares $\ket{\psi_k}=C_k\ket{\psi_0}$. At each iteration, we record a collection of $N$ independent classical shadows obtained from uniformly random local Pauli measurements of $\ket{\psi_k}$, which allows us to estimate the energy difference
\begin{equation}
   \Delta E_{k, j}(\theta) =\bra{\psi_k}H\ket{\psi_k}-\bra{\psi_k}U^\dagger _j(\theta)HU_j(\theta)\ket{\psi_k},
\end{equation}
for every candidate gate $U_j(\theta)$ in a given set $\mathcal{G}=\{U_j(\theta)\}^K_{j=1}$. Classically maximizing the estimated energy difference over $\theta$ for all candidate gates enables us to identify the energy-minimizing gate $U_{j^*}(\theta^*_{j^*})$ and update the circuit as $C_k\to C_{k+1}=U_{j^*}(\theta^*_{j^*})C_k$. The procedure is repeated until either a predefined circuit depth $D$ is reached or no candidate gate yields an energy reduction exceeding a predefined threshold $\Delta >0$. 

To understand how SEGQE estimates the energy-minimizing parameters $\theta^*_j$ for each gate in $\mathcal{G}$ at iteration $k$ (for notational convenience, we drop the index $k$ in the following), consider a single arbitrarily parametrized $n$-qubit unitary $U(\theta) \in \mathrm{SU}(2^n)$ of locality $m\leq n$. 
Denote by $Q^u\subseteq\{1,\ldots n\}$ the qubit support of $U(\theta)$, and by $Q^c$ its complement, on which $U(\theta)$ acts trivially. Likewise, let $Q^i\subseteq\{1,\ldots n\}$ denote the support of the Hamiltonian term $P_i$. We can decompose each Pauli $P_i$ with respect to the bipartition $Q^u\cup Q^c$ as $P_i=P^u_i\otimes P^c_i$, where $P^u_i$ acts on $Q^u$ and $P^c_i$ acts on $Q^c$. Furthermore, we define the set $I^u=\left\{i\in\{1,\ldots,r\}: Q^i\cap Q^u\neq \emptyset\right\}$ as the indices of the Hamiltonian terms which act nontrivially on the qubits in $Q^u$. Hamiltonian terms with $Q^i \cap Q^u = \emptyset$ commute with $U(\theta)$ and therefore do not contribute to the energy difference. The energy difference induced by $U(\theta)$ is then given by
\begin{align}
    \Delta E(\theta)=&\bra{\psi}H\ket{\psi}-\bra{\psi}U^\dagger(\theta) H U(\theta) \ket{\psi}\\
	=& \sum_{i\in I^u}c_i 
	\left(\bra{\psi}P_i\ket{\psi}{-}\bra{\psi}U^\dagger(\theta) P^u_i U(\theta)\otimes P_i^c\ket{\psi} \right) \nonumber
\end{align}
    Since the $m$-qubit Pauli operators $\{P_\mu^u\}_{\mu=0}^{4^m-1}$ form an orthonormal basis (with respect to the Hilbert--Schmidt inner product) on the operator space corresponding to the qubits in $Q^u$, we can expand
    \begin{align}
        U^\dagger(\theta) P^u_i U(\theta) &= \sum_{\mu=0}^{4^m-1}\frac{1}{2^m}\mathrm{tr}\left[U^\dagger(\theta) P^u_iU(\theta)P_\mu^u\right]P_\mu^u\\&=\sum_{\mu=0}^{4^m-1}r_{i\mu}P_\mu^u,
    \end{align}
    where $r_{i\mu}(\theta):=\frac{1}{2^m}\mathrm{tr}\left[U^\dagger(\theta) P^u_iU(\theta)P_\mu^u\right]$. Substituting this into the above expression for the energy difference yields
    \begin{equation}
        \label{eq:energy_differences}
        \Delta E(\theta)=\sum_{i\in I^u}\sum^{4^m-1}_{\mu=0}f_{i\mu}(\theta)p_{i\mu},
    \end{equation}
    where we define
    \begin{equation*}
        f_{i\mu}(\theta):= c_i(\delta_{P^u_i,P_\mu^u} - r_{i\mu}(\theta)),\quad p_{i\mu} := \bra{\psi} P_\mu^u \otimes P_i^c \ket{\psi}.
    \end{equation*}
    Notably, the expectation values $p_{i\mu}$ are independent of $\theta$, so once estimated, $\Delta E(\theta)$ can be constructed via \cref{eq:energy_differences} and thus evaluated and maximized entirely classically. To simplify notation, we define the composite index $\alpha:=(i,\mu)$. For each candidate gate $U_j(\theta)\in\mathcal{G}$, we define the associated Pauli operators $P_{j,\alpha}:=P^{u(j)}_\mu \otimes P^{c(j)}_i$, and their expectation values $p_{j,\alpha}:=\bra{\psi}P_{j,\alpha}\ket{\psi}$, where the superscripts $u(j)$ and $c(j)$ indicate the bipartition induced by the support of $U_j$. The corresponding energy differences can then be written as
    \begin{equation}
        \Delta E_j(\theta) =\sum_{\alpha\in \mathcal{F}_j}f_{j,\alpha}(\theta)p_{j,\alpha}, 
    \end{equation}
    where $\mathcal{F}_j$ denotes the set of indices $\alpha$ that appear in the expansion associated with gate $U_j$. 
    In the following section, we derive rigorous guarantees on the estimation accuracy achievable when using classical shadows to predict these energy differences. The overall procedure of SEGQE is summarized in \cref{alg:segqe}.

\begin{figure}[htbp]
    \begin{minipage}{1\linewidth}
        \begin{algorithm}[H]
  \caption{Shadow Enhanced Greedy Quantum Eigensolver (SEGQE)}
  \label{alg:segqe}
  \begin{algorithmic}[1]
    \Require Hamiltonian $H = \sum_{i=1}^r c_i P_i$, initial state $\ket{\psi_0}$,
             gate set $\mathcal{G} = \{U_j(\theta)\}_{j=1}^K$, 
             maximum depth $D$, threshold $\Delta > 0$
    \State $C_0 \gets I$, $k \gets 0$
    \While{$k < D$}
      \State $\ket{\psi_k} \gets C_k \ket{\psi_0}$
      \State Record $N$ independent shadows $\{\hat{\rho}_k^{(t)}\}_{t=1}^N$ of $\ket{\psi_k}$
      \State Use shadows to estimate all required $p_{j,\alpha}$ 
      \State $\Delta E_{\mathrm{max}} \gets 0$, $(j^*, \theta^*) \gets \text{None}$
      \ForAll{candidate gates $U_j(\theta) \in \mathcal{G}$}
        \State Use $p_{j,\alpha}$ to classically find $\theta^*_j=\underset{\theta}{\arg\max} \Delta E_{k,j}(\theta)$
        \If{$\Delta E_{k,j}(\theta^*_j) > \Delta E_{\mathrm{max}}$}
          \State $\Delta E_{\mathrm{max}} \gets \Delta E_{k,j}(\theta^*_j)$
          \State $(j^*, \theta^*) \gets (j, \theta^*_j)$
        \EndIf
      \EndFor
      \If{$\Delta E_{\mathrm{max}} \le \Delta$}
        \State \textbf{break}  %\Comment{no sufficiently large energy reduction}
      \EndIf
      \State $C_{k+1} \gets U_{j^*}(\theta^*) C_k$
      \State $k \gets k + 1$
    \EndWhile
    \State \Return $C_k$%, $E_k$
  \end{algorithmic}
\end{algorithm}
    \end{minipage}

\end{figure}




